\section{矩阵函数：从多项式到指数}
\label{sec:03-Linear-Algebra/05-Eigenvalues-and-Eigenvectors/04}

前两节各自算出了一个东西：可对角化时 \(A^k=S\Lambda^kS^{-1}\)，有 Jordan 块时 \(e^{Jt}\) 带出 \(t\) 的幂次。两次的手法看着像是各自碰巧凑效。把它们放到一起，该问的问题就浮出来了——\(e^A\)、\(\sqrt{A}\)、\(\log A\) 到底是什么意思，凭什么这么写。

先排除一个诱人的错答案。既然矩阵是一堆数，逐个元素代入函数不就完了？拿上一节那个幂零矩阵试：

\[
A=\begin{pmatrix}0&1\\0&0\end{pmatrix},
\qquad
\bigl(e^{a_{ij}}\bigr)=\begin{pmatrix}1&e\\1&1\end{pmatrix},
\qquad
e^{A}=I+A=\begin{pmatrix}1&1\\0&1\end{pmatrix}.
\]

右边那个是真的——因为 \(A^2=0\)，幂级数只剩两项。两个结果连形状都不一样，逐元素的那个甚至不再是上三角，谈不上和 \(A\) 有什么关系。逐元素运算把矩阵当成了数组；矩阵函数要尊重的是它作为线性映射的结构。

\subsection{多项式是无争议的起点}

只用加法和乘法就能定义的那部分，不会有歧义。给定

\[
p(z)=c_0+c_1z+\cdots+c_mz^m,
\qquad
p(A)=c_0I+c_1A+\cdots+c_mA^m .
\]

常数项必须写成 \(c_0I\)，因为矩阵乘法的单位元是 \(I\) 而不是数 \(1\)。整个表达式只涉及同一个 \(A\) 的幂，它们彼此交换，因此像标量多项式那样合并同类项、做因式分解都合法。

有意思的是矩阵会让高次项自己折回来。上面那个 \(A\) 满足 \(A^2=0\)，所以无论多项式次数多高，\(p(A)=c_0I+c_1A\)。一般情形由 Cayley--Hamilton 定理管着：方阵满足自己的特征多项式，\(\chi_A(A)=0\)，于是任何 \(p\) 都可以对 \(\chi_A\) 取余，化成次数低于 \(n\) 的多项式。想更紧凑就用上一节的最小多项式：两个多项式只要差一个 \(m_A\) 的倍数，作用在 \(A\) 上的结果就完全相同。

这条观察给出一个判断：\(f(A)\) 能取到的值不超过 \(n\) 个自由度，它由 \(f\) 在谱上的行为决定，跟 \(f\) 在别处长什么样无关。

\subsection{可对角化时，函数只作用在特征值上}

设 \(A=S\Lambda S^{-1}\)。代入多项式，中间的 \(S^{-1}S\) 成对消掉，

\[
p(A)=S\,p(\Lambda)\,S^{-1},
\qquad
p(\Lambda)=\diag\bigl(p(\lambda_1),\dots,p(\lambda_n)\bigr).
\]

推广的方向就此明确：只要标量函数 \(f\) 在所有特征值上有定义，就令

\[
f(A)=S\,\diag\bigl(f(\lambda_1),\dots,f(\lambda_n)\bigr)\,S^{-1}.
\]

这个定义的意思用动力学的话说最清楚：\(A\) 把第 \(i\) 个特征模态缩放 \(\lambda_i\) 倍，\(f(A)\) 就把它缩放 \(f(\lambda_i)\) 倍，不变方向一条也没变，变的只是每条方向上的那个数。定义还不依赖特征向量的长度和排列——改变长度会在 \(S\) 与 \(S^{-1}\) 里抵消，交换列会同时交换对角元。

拿本单元反复用的 \(A=\begin{pmatrix}2&1\\1&2\end{pmatrix}\) 算个平方根。它对称，特征方向正交，取

\[
Q=\frac1{\sqrt2}\begin{pmatrix}1&1\\1&-1\end{pmatrix},
\qquad
A=Q\diag(3,1)Q^{\trans},
\]

于是 \(A^{1/2}=Q\diag(\sqrt3,1)Q^{\trans}=\frac12\begin{pmatrix}\sqrt3+1&\sqrt3-1\\\sqrt3-1&\sqrt3+1\end{pmatrix}\)，乘开可验证平方回到 \(A\)。这里悄悄做了一次选择：每个特征值都取了正的标量根。换一组符号会得到另一个平方根，允许复数则更多，某些实矩阵甚至没有实平方根。记号 \(A^{1/2}\) 只有在说清分支之后才是确定的对象；对称正定矩阵最省事，正特征值加正交对角化，结果唯一且仍是对称正定的。

\subsection{Jordan 块还要读函数的导数}

不可对角化时，只知道 \(f\) 在特征值上的取值不够。对大小为 \(r\) 的块 \(J=\lambda I+N\)（\(N^r=0\)），把标量的 Taylor 展开中的增量换成 \(N\)：

\[
f(J)=f(\lambda)I+f'(\lambda)N+\frac{f''(\lambda)}{2!}N^2+\cdots+\frac{f^{(r-1)}(\lambda)}{(r-1)!}N^{r-1},
\]

级数在 \(r-1\) 次后自动断掉。若 \(A=SJS^{-1}\)，就逐块用这条式子再拼回去。最小的二阶块已经把话说完：

\[
J=\begin{pmatrix}\lambda&1\\0&\lambda\end{pmatrix}
\ \Longrightarrow\
f(J)=\begin{pmatrix}f(\lambda)&f'(\lambda)\\0&f(\lambda)\end{pmatrix}.
\]

对角线读函数值，超对角线读一阶导数，更长的链继续读更高阶。缺陷结构记住的不只是谱在哪里，还有函数在谱附近怎么变。这也划出了定义的光滑性要求：最大块为 \(r\) 时，\(f\) 至少要在相应特征值处有到 \(r-1\) 阶的导数；只在特征值上给了一堆数值而没有导数信息的``函数''，在缺陷矩阵上根本定不出唯一的 \(f(A)\)。

\subsection{矩阵指数把离散步补成连续轨迹}

指数函数处处解析，用幂级数 \(e^{tA}=\sum_k t^kA^k/k!\) 直接定义即可，前面几节已经用过。它最值得看的一面是复特征值对应的情形。取

\[
A=\begin{pmatrix}\alpha&-\omega\\\omega&\alpha\end{pmatrix}=\alpha I+\omega J,
\qquad
J=\begin{pmatrix}0&-1\\1&0\end{pmatrix},\quad J^2=-I .
\]

\(\alpha I\) 与 \(\omega J\) 交换，指数可以拆开；把 \(e^{\omega tJ}\) 的级数按奇偶次分组，用 \(J^{2k}=(-1)^kI\) 与 \(J^{2k+1}=(-1)^kJ\)，两组恰好凑成余弦与正弦：

\[
e^{tA}=e^{\alpha t}\begin{pmatrix}\cos\omega t&-\sin\omega t\\\sin\omega t&\cos\omega t\end{pmatrix}.
\]

第一节说复特征值 \(\alpha\pm i\omega\) 意味着旋转，这就是那句话的显式版本：实部管半径的胀缩，虚部管转速。

\begin{center}
\begin{tikzpicture}[x=1.15cm,y=1.15cm,>=stealth,font=\small]
  \draw[gray!55,->] (-2.2,0) -- (2.4,0);
  \draw[gray!55,->] (0,-1.7) -- (0,2.0);
  \draw[blue!70!black,thick,domain=0:6,samples=160]
    plot ({2*exp(-0.15*\x)*cos(60*\x)},{2*exp(-0.15*\x)*sin(60*\x)});
  \draw[black!45,dashed]
    (2,0) -- (0.861,1.491) -- (-0.741,1.283) -- (-1.276,0)
    -- (-0.549,-0.951) -- (0.472,-0.818) -- (0.813,0);
  \foreach \k in {0,1,2,3,4,5,6}
    \fill[black] ({2*exp(-0.15*\k)*cos(60*\k)},{2*exp(-0.15*\k)*sin(60*\k)}) circle (1.8pt);
  \node[below right] at (2,-0.05) {\(x_0\)};
  \node[above right] at (0.9,1.5) {\(x_1\)};
  \node[align=center] at (-2.0,-1.35) {实线：\(e^{tA}x_0\)\\虚线：整数步的迭代};
\end{tikzpicture}

\emph{一个 \(\operatorname{Re}\lambda<0\) 的复特征值对：轨迹旋进原点。离散迭代只看得到黑点，矩阵指数把它们之间的路补齐；反过来，\(A_{\mathrm d}=e^{hA}\) 正是采样间隔为 \(h\) 时的一步转移矩阵。}
\end{center}

图里那句反过来的话在工程上用得最多。连续系统 \(\dot x=Ax+Bu\) 的解由常数变易公式给出，

\[
x(t)=e^{tA}x(0)+\int_0^t e^{(t-s)A}Bu(s)\,\mathrm ds,
\]

若输入在每个采样周期内保持不变，把 \(t\) 取成 \(h\) 就得到一个精确的离散模型：状态矩阵是 \(e^{hA}\)，输入矩阵是 \(\bigl(\int_0^h e^{sA}\mathrm ds\bigr)B\)。控制器在离散时间里设计、在连续对象上运行，靠的就是这次换算，而稳定性判据也随之从 \(\operatorname{Re}\lambda<0\) 变成 \(\abs{e^{h\lambda}}<1\)——两者其实等价，因为指数把左半平面映成单位圆内部。类似的换算还出现在连续时间 Markov 链里：生成元 \(Q\) 与转移矩阵的关系是 \(P(t)=e^{tQ}\)，\hyperref[ch:050]{``连续时间 Markov 链与排队：不仅要知道跳到哪，还要知道何时跳''}会用到它。近年把网络层数当作连续时间的做法（神经常微分方程、扩散模型的连续极限）也是同一个思路：先写下 \(\dot x=f(x,t)\)，线性化之后主导行为仍由 \(e^{tA}\) 描述。

若特征值带 Jordan 块，\(e^{tA}\) 里就会出现 \(t^je^{\lambda t}\)，二阶块给出 \(e^{\lambda t}(I+tN)\)——上一节的多项式增长，在这里表现为函数值旁边多了一个导数项。

\subsection{标量恒等式不能照抄}

\(e^{x+y}=e^xe^y\) 对数是天经地义的，对矩阵只有在 \(AB=BA\) 时才成立。展开到二次项就能看出分歧：

\[
\begin{aligned}
e^{A+B}&=I+A+B+\tfrac12\bigl(A^2+AB+BA+B^2\bigr)+\cdots,\\
e^Ae^B&=I+A+B+\tfrac12A^2+AB+\tfrac12B^2+\cdots,
\end{aligned}
\]

差别落在交叉项上，差额由对易子 \(AB-BA\) 度量。这不是纸面上的洁癖：两个不同的控制作用先后施加与调换顺序，结果确实不同，机器人靠交替动作实现``侧向平移''就是这个非交换性的直接后果。数值上则有 Trotter 分裂，把难算的 \(e^{h(A+B)}\) 近似成 \(e^{hA}e^{hB}\)，代价是每步 \(O(h^2)\) 的误差，正好来自上式的差额。同理，\(\log(AB)=\log A+\log B\) 与 \((AB)^t=A^tB^t\) 都需要额外条件，不能凭标量直觉搬用。

另一条边界来自谱。要得到实的对数或实的平方根，仅仅``标量函数在正数上有定义''是不够的：负实特征值、缺陷结构、复分支的选择都会影响存在性与唯一性。对称正定矩阵仍然是最规矩的一类，\(f\) 逐个作用在正特征值上，结果保持实对称——\hyperref[ch:022]{``对称矩阵、二次型与正定性''}会把这类矩阵的好性质集中处理。

\subsection{定义方式与计算方式是两码事}

Jordan 形给出了最统一的定义公式，却因为上一节讲过的扰动敏感性，不适合直接拿来做浮点算法。实际计算 \(e^{A}\) 的主流办法是 scaling and squaring：先把 \(A\) 除以一个 \(2^m\) 使范数变小，用 Pad\'e 有理逼近算出 \(e^{A/2^m}\)，再平方 \(m\) 次还原。需要 Schur 分解时，把矩阵化成上三角后逐块递推也很稳。而在大规模问题里，人们往往并不需要整个 \(e^{A}\)，只需要它作用在某个向量上的结果 \(e^{A}v\)；这时 Krylov 子空间方法可以在不形成稠密矩阵的前提下给出高精度近似，这也是求解大型线性演化方程的常规做法。

由此收束一句可以带走的判断：知道一个对象怎样定义，不等于知道该怎样算它。定义要求普适与唯一，算法要求在有限精度下不放大误差，两种要求经常指向不同的表示。

本单元到此把``反复作用一个矩阵''这件事拆完了：找不变方向，换到那组方向上，让每条方向独立地取幂或取指数；方向不够时补上链节，代价是多出多项式因子。整套机器最顺手的场合，是特征方向恰好互相垂直的时候——那时 \(S^{-1}\) 就是 \(S^{\trans}\)，一切数值上的隐患随之消失。下一单元先给出这种情形的判据，再看它怎样把二次型、正定性与曲率联系起来。
